% Generated draft: requires mathematical and source review.
\documentclass[11pt]{article}
\usepackage[margin=25mm]{geometry}
\usepackage{fontspec}
\setmainfont{Latin Modern Roman}
\usepackage{amsmath,amssymb,mathtools,unicode-math}
\setmathfont{Latin Modern Math}
\usepackage{enumitem,needspace,xurl,hyperref}
\setlength{\parindent}{0pt}
\setlength{\parskip}{5pt}
\setlength{\emergencystretch}{3em}
\newcommand{\N}{\mathbb{N}}
\newcommand{\Z}{\mathbb{Z}}
\newcommand{\Q}{\mathbb{Q}}
\newcommand{\R}{\mathbb{R}}
\newcommand{\C}{\mathbb{C}}
\newcommand{\F}{\mathbb{F}}
\newcommand{\A}{\mathbb{A}}
\newcommand{\PP}{\mathbb{P}}
\newcommand{\E}{\mathbb{E}}
\newcommand{\eps}{\varepsilon}
\newcommand{\dd}{\,\mathrm{d}}
\begin{document}
\section*{UnsolvedMath ID: 700836}
\textbf{Source problem code:} catalogue entry
\subsection{Definitions and mathematical statement}
For a model \(M\), an automorphism is a bijection of its underlying set that preserves all interpretations in its language. Once an effective presentation of the underlying set has been specified, call an automorphism recursive if its action is a total computable function in that presentation. Under the interpretation that the source means the full group of such automorphisms, write
\[
\operatorname{Aut}_{\mathrm{rec}}(M)
=\{f\in\operatorname{Aut}(M):f\text{ is recursive in the specified presentation}\}.
\]
Let \(\mathcal C\) consist of the groups isomorphic to \(\operatorname{Aut}_{\mathrm{rec}}(M)\) for models \(M\) allowed by the intended presentation convention. No computability assumption on the interpretations of the language is added here.

For groups \(G,H\), their free product \(G*H\) is the group with homomorphisms \(\iota_G:G\to G*H\) and \(\iota_H:H\to G*H\) such that, for every group \(K\) and homomorphisms \(u:G\to K\), \(v:H\to K\), there is a unique homomorphism \(w:G*H\to K\) satisfying \(w\circ\iota_G=u\) and \(w\circ\iota_H=v\).

\textbf{Question (binary free-product interpretation).} Is it true that
\[
\forall G,H\in\mathcal C,\qquad G*H\in\mathcal C?
\]
Equivalently, for every pair of models \(M,N\) in the intended framework, does there exist a model \(P\) in that framework such that
\[
\operatorname{Aut}_{\mathrm{rec}}(P)
\cong
\operatorname{Aut}_{\mathrm{rec}}(M)*\operatorname{Aut}_{\mathrm{rec}}(N)?
\]
\textbf{Unresolved qualifications.} The supplied statement refers to ``arbitrary models'' but specifies neither their languages nor their underlying domains or effective presentations. Thus the admissible models and the meaning of recursiveness remain parameters of this formulation. The full-group interpretation above also requires confirmation. Furthermore, the source does not explicitly distinguish binary free products from free products of arbitrary indexed families; the latter would ask whether \(\mathop{*}_{i\in I}G_i\in\mathcal C\) for every allowed index set \(I\) and every family \((G_i)_{i\in I}\) in \(\mathcal C\). The permitted index sets are not supplied.
\subsection{Short English statement}
S. S. Goncharov asks whether groups arising as groups of recursive automorphisms of models are closed under free products. Under the full-group and binary interpretations, any two such groups should have a free product that is again the full recursive automorphism group of some model. The effective-presentation conventions and the intended scope of free products need clarification.
\subsection{Sources}
\begin{itemize}
\item Local UnsolvedMath registry: permanent ID \(700836\), exact source problem code \texttt{catalogue entry}, title ``Kourovka Notebook Problem 12.16.'' This draft uses the supplied complete source record, with upstream revision \texttt{d371ee917ac9647fbb06cd6fc529deec6f3dda82}.
\item Supplied attribution: S. S. Goncharov, proposed in 1992; \emph{The Kourovka Notebook (Unsolved Problems in Group Theory)}, 12th issue (1992), Problem 12.16.
\item The supplied background also identifies the 21st issue edition (2026), \texttt{arXiv:1401.0300v48}. The supplied collection description attributes the Notebook's editing to E. I. Khukhro and V. D. Mazurov. These bibliographic details have not been independently checked.
\item The registry labels the problem \texttt{open}; that label is supplied source metadata, not an independently verified assessment of its mathematical or literature status.
\end{itemize}
\end{document}
